In this section, we provide a brief overview of the language features of Rust
that commonly appear in the soundness issues we collected.

\subsection{Lifetimes}
\label{sec:background:lifetime}

In Rust, a \emph{reference} is a pointer with a static guarantee that it points
to a valid value.
%
Each reference has a \emph{lifetime}, which specifies the duration for which the
reference is valid.
%
Lifetimes are key to ensure memory safety;
%
for instance, the compiler ensures that a lifetime ends before the referenced
value goes out of scope.
%
A reference has a type \C{\&'a T}, where \C{'a} is the lifetime and \C{T} is the
type of the value.

Functions can have \emph{lifetime parameters} to express the relationship
between the lifetimes of the parameters and the return value.
%
For example, the following function with two lifetime parameters \C{'a} and
\C{'b} takes two references to integers (\C{i32}) and returns the first one,
specifying that the lifetime of the return value is the same as that of the
first parameter:

\vspace{-0.5em}
\begin{lstlisting}[basicstyle=\linespread{0.85}\ttfamily\footnotesize,]
fn foo<'a, 'b>(v: &'a i32, u: &'b i32) -> &'a i32 { v }
\end{lstlisting}
\vspace{-0.5em}

When a lifetime \C{'a} is longer than or equal to another lifetime \C{'b}, i.e.,
\C{'a} \emph{outlives} \C{'b}, a reference of lifetime \C{'a} can be used where
a reference of lifetime \C{'b} is expected.
%
By default, different lifetime parameters are not considered to outlive each
other, but programmers can introduce \emph{lifetime bounds} of the form \C{'a:
'b}, implying \C{'a} outlives \C{'b}, as shown below:

\vspace{-0.5em}
\begin{lstlisting}[basicstyle=\linespread{0.85}\ttfamily\footnotesize,]
fn foo<'a, 'b>(v: &'a i32) -> &'b i32 { v } // error
fn bar<'a: 'b, 'b>(v: &'a i32) -> &'b i32 { v } // ok
fn baz() { let p; { let v = 0; p = bar(&v); } *p; } // error
fn qux() { let v = 0; let p = bar(&v); *p; } // ok
\end{lstlisting}
\vspace{-0.5em}

\noindent
%
In \C{foo}, returning \C{v} is disallowed because there is no guarantee that
\C{'a} outlives \C{'b};
%
in \C{bar}, \C{'a: 'b} allows it.
%
The compiler checks whether the specified lifetime bounds are satisfied at each
call site:
%
\C{baz} is rejected because \C{p} is used after \C{v} goes out of scope, while
\C{qux} is accepted.

%Compound data types (structs, enums, and unions) can also have lifetime
%parameters to specify the lifetimes of references contained within them, as in
%the following example:

%\vspace{-0.5em}
%\begin{lstlisting}
%struct S<'a> { x: i32, y: &'a i32 }
%\end{lstlisting}
%\vspace{-0.5em}

%\noindent
%%
%Here, \C{S} has a lifetime parameter \C{'a} for the reference \C{y}.
%%
%Since \C{y} is valid only within \C{'a}, a value of type \C{S<'a>} is also valid
%only within \C{'a}.
%%
%This naturally extends the notion of lifetime bounds to types:
%%
%a type \C{T} outlives \C{'a}, written \C{T: 'a}, if all lifetimes occurring in
%\C{T} outlive \C{'a}.
%%
%For example, when \C{'a: 'b}, both \C{\&'a i32: 'b} and \C{S<'a>: 'b} hold.

Since a value must remain valid while its reference is live, a type \C{\&'a T}
is considered \emph{well-formed} only when every lifetime appearing in \C{T}
outlives \C{'a}.
%
For example, \C{\&'a \&'b i32} is well-formed only if \C{'b: 'a} holds.

To reduce the burden of stating lifetime bounds, Rust supports \emph{implied
bounds}, which automatically introduce lifetime bounds deduced from the
well-formedness of the types used in each definition (e.g., the parameter and
return types of a function).
%
In the following example, \C{foo} passes type checking due to the implied bound
\C{'b: 'a} and is equivalent to \C{bar}:

\vspace{-0.5em}
\begin{lstlisting}[basicstyle=\linespread{0.85}\ttfamily\footnotesize,]
fn foo<'a, 'b>(v: &'a &'b i32) {}
fn bar<'a, 'b: 'a>(v: &'a &'b i32) {}
\end{lstlisting}
\vspace{-0.5em}

\noindent
%
Note that implied bounds are distinct from Rust's lifetime elision, which allows
programmers to omit explicit lifetime annotations (e.g., writing \C{\&i32}
instead of \C{\&'a i32}), not lifetime bounds, which the compiler then infers
based on fixed syntactic patterns.

So far, the names of lifetimes have been explicitly introduced by lifetime
parameters, but there is one exception:
%
the special lifetime \C{'static} denotes the entire duration of the program and
does not need to be defined as a lifetime parameter.
%
By definition, \C{'static} outlives any lifetime.
%
A reference to a global variable or a constant has the \C{'static} lifetime.

% \vspace{-0.5em}
% \begin{lstlisting}
% static X: i32 = 0;
% fn foo() { let x: &'static i32 = &X; let y: &'static i32 = &0; }
% \end{lstlisting}
% \vspace{-0.5em}

\subsection{Traits}
\label{sec:background:trait}

A \emph{trait} represents a set of types, similar to abstract classes or
interfaces in object-oriented languages and type classes in Haskell.
%
Consider the following example:

\vspace{-0.5em}
\begin{lstlisting}[basicstyle=\linespread{0.85}\ttfamily\footnotesize,]
trait X { fn foo(self); } $\label{line:trait:1}$
struct S; impl X for S { fn foo(self) { println!("S"); } } $\label{line:trait:2}$
fn bar<T: X>(v: T) { v.foo(); } $\label{line:trait:3}$
fn baz<T>(v: T) where T: X { v.foo(); } $\label{line:trait:4}$
fn qux(v: impl X) { v.foo(); } $\label{line:trait:5}$
\end{lstlisting}
\vspace{-0.5em}

\noindent
%
Here, \C{X} denotes a set of types that implement the method \C{foo}
(line~\ref{line:trait:1}), and \C{S} belongs to \C{X}, i.e., \C{S}
\emph{implements} \C{X} (line~\ref{line:trait:2}).
%
The function \C{bar} takes a value of any type \C{T} that implements \C{X} and
calls \C{foo} (line~\ref{line:trait:3}).
%
In \C{bar}, \C{X} is a \emph{trait bound} of \C{T}.
%
Trait bounds can also be specified in a \C{where} clause
(line~\ref{line:trait:4}).
%
In addition, instead of introducing a type parameter and its bound, we can use
an \emph{impl trait} type, such as \C{impl X}, as a parameter type
(line~\ref{line:trait:5}).
%
This use of an \C{impl} trait type is called an \emph{argument-position impl
trait} (APIT).
%
All the functions \C{bar}, \C{baz}, and \C{qux} are equivalent, representing the
same set of possible arguments with different syntax.
%
% Note that a type can have multiple bounds, separated by the \C{+} symbol
% (e.g., \C{X + Y} or \C{X + 'a}).

It is possible to implement a trait for multiple types using a single \C{impl},
as shown below:

\vspace{-0.5em}
\begin{lstlisting}
struct S; struct R; struct Q; trait X {} impl X for S {} impl X for R {} $\label{line:blanket1}$
trait Y {} impl<T: X> Y for T {} trait Z {} impl<T> Z for T {} $\label{line:blanket2}$
\end{lstlisting}
\vspace{-0.5em}

\noindent
%
\C{S} and \C{R} implement \C{X}, but \C{Q} does not (line~\ref{line:blanket1}).
%
Since \C{Y} is implemented for all types that implement \C{X}, both \C{S} and
\C{R} implement \C{Y}, but \C{Q} does not;
%
on the other hand, \C{Z} is implemented for all types, including \C{S}, \C{R},
and \C{Q} (line~\ref{line:blanket2}).

Traits can require types not only to provide particular methods but also to
specify \emph{associated types}, as in the following example:

\vspace{-0.5em}
\begin{lstlisting}
trait X { type A; fn foo(self) -> Self::A; } $\label{line:assoc:1}$
struct S; impl X for S { type A = i32; fn foo(self) -> Self::A { 0 } } $\label{line:assoc:2}$
struct R; impl X for R { type A = f32; fn foo(self) -> Self::A { 0.0 } } $\label{line:assoc:3}$
fn bar<T: X>(v: T) -> T::A { v.foo() } $\label{line:assoc:4}$
fn baz<T: X>(v: T) -> <T as X>::A { v.foo() } $\label{line:assoc:5}$
fn qux<T: X<A = i32>>(v: T) -> i32 { v.foo() } $\label{line:assoc:6}$
\end{lstlisting}
\vspace{-0.5em}

\noindent
%
A type implementing \C{X} must specify an associated type \C{A}, and its method
\C{foo} must return a value of type \C{A} (line~\ref{line:assoc:1}).
%
Both \C{S} and \C{R} implement \C{X}, defining different types for \C{A}
(lines~\ref{line:assoc:2} and \ref{line:assoc:3}).
%
The function \C{bar} returns the value from \C{foo}, whose type is denoted
\C{T::A} (line~\ref{line:assoc:4}).
%
We can also write \C{<T as X>::A} to disambiguate which trait declares the
associated type (line~\ref{line:assoc:5}), which is useful when multiple traits
define an associated type with the same name.
%
Additionally, to specify a set of types implementing a trait with a specific
associated type, we can use a trait bound like \C{X<A = i32>};
%
this allows \C{S} to be passed to \C{qux}, but not \C{R}
(line~\ref{line:assoc:6}).

Traits are also used to hide the concrete return type of a function by
specifying the return type as \C{impl X}, called a \emph{return-position impl
trait} (RPIT).
%
Consider the example below:

\vspace{-0.5em}
\begin{lstlisting}
trait X { fn foo(self); } $\label{line:rpit:1}$
struct S; impl X for S { fn foo(self) {} } impl S { fn bar(self) {} } $\label{line:rpit:2}$
fn baz() -> impl X { S } $\label{line:rpit:3}$
fn qux() { let v = baz(); v.foo(); /* ok */ v.bar(); /* error */ } $\label{line:rpit:4}$
\end{lstlisting}
\vspace{-0.5em}

\noindent
%
\C{S} defines two methods \C{foo} and \C{bar}, while only \C{foo} is part of the
trait \C{X} (lines~\ref{line:rpit:1} and \ref{line:rpit:2}).
%
The function \C{baz} returns \C{S}, but it hides this information from callers,
revealing only that the return type implements \C{X} (line~\ref{line:rpit:3}).
%
Therefore, \C{qux} can call \C{foo} on the returned value, but not \C{bar}
(line~\ref{line:rpit:4}).
%
Note that the purpose of an RPIT is not to allow returning values of different
types.
%
Even when using an RPIT, a single underlying return type must exist.
%
For instance, the following code is rejected even though both \C{S} and \C{R}
implement \C{X}:

\vspace{-0.5em}
\begin{lstlisting}
trait X {} struct S; impl X for S {} struct R; impl X for R {}
fn foo(v: i32) -> impl X { if v == 0 { S } else { R } } // error
\end{lstlisting}
\vspace{-0.5em}

To enable returning values of different types, Rust supports \emph{trait
objects}, which accompany values with virtual tables for dynamic dispatch.
%
The type of a trait object is \C{dyn X}, where \C{X} is the implemented trait,
as in the following example:

\vspace{-0.5em}
\begin{lstlisting}
fn foo(v: i32) -> &'static dyn X { if v == 0 { &S } else { &R } } // ok
\end{lstlisting}
\vspace{-0.5em}

\noindent
%
Note that a trait object is \emph{dynamically sized}, i.e., its size is unknown
at compile time, because values of different types can be stored.
%
It is thus typically placed behind a pointer.

\subsection{Higher-Rank Polymorphism}
\label{sec:background:hr}

Rust supports rank-1 type polymorphism (e.g., generic functions), but not
higher-rank type polymorphism (e.g., functions that take generic functions).
%
However, it supports higher-rank \emph{lifetime} polymorphism because it is
extremely useful when dealing with function pointers:

\vspace{-0.5em}
\begin{lstlisting}
fn foo<'a>(v: &'a i32) -> i32 { *v }$\label{line:hr:1}$
fn bar<'a>(f: fn(&'a i32)->i32) {{let v=0;f(&v);} {let u=1;f(&u);}} // error$\label{line:hr:2}$
fn baz(f: for<'a> fn(&'a i32)->i32) {{let v=0;f(&v);} {let u=1;f(&u);}}// ok$\label{line:hr:3}$
\end{lstlisting}
\vspace{-0.5em}

\noindent
%
The function \C{foo} takes a reference of any lifetime (line~\ref{line:hr:1}).
%
Another function may want to take a function pointer to \C{foo} and call it on
references with different lifetimes.
%
This cannot be achieved with rank-1 polymorphism:
%
\C{bar} takes a function that accepts a reference of \emph{some} lifetime,
failing to pass both \C{\&v} and \C{\&u} (line~\ref{line:hr:2}).
%
With higher-rank polymorphism, \C{baz} takes a function that accepts a reference
of \emph{any} lifetime, allowing both calls (line~\ref{line:hr:3}).

Higher-rank polymorphism is also useful when reference types implement traits:

\vspace{-0.5em}
\begin{lstlisting}
trait X { fn foo(self); }$\label{line:hrtb:1}$
fn bar<T>(mut v: T) where for<'a> &'a T: X {$\label{line:hrtb:2}$
  { let u = v; (&u).foo(); v = u; } { let w = v; (&w).foo(); } }$\label{line:hrtb:3}$
struct S; impl<'a> X for &'a S { fn foo(self) {} } fn baz() { bar(S); }$\label{line:hrtb:4}$
\end{lstlisting}
\vspace{-0.5em}

\noindent
%
Using a \emph{higher-rank trait bound} (HRTB), the function \C{bar} requires the
type parameter \C{T} to satisfy the condition that \C{\&'a T} implements \C{X}
for any lifetime \C{'a} (line~\ref{line:hrtb:2}).
%
This allows the function body to call the method \C{foo}, which is provided by \C{X},
using references of different lifetimes (line~\ref{line:hrtb:3}).
%
Since \C{\&'a S} implements \C{X} for any lifetime \C{'a}, \C{baz} can pass
\C{S} to \C{bar} (line~\ref{line:hrtb:4}).
